\begin{table}[htbp]
  \centering
  \caption{Effect of the allocation ratio on the required sample size, under within-group independence between response and dropout. Fixed parameters: $\alpha = 0.025$ (one-sided), target power $= 0.8$.}
  \label{tab:allocation}
  \begin{threeparttable}
  \begin{tabular}{ccccrrr}
    \toprule
    $\pi_{1}$ & $\pi_{0}$ & $\omega$ & $r$ & $N_{\mathrm{CC}}$ & $N_{\mathrm{NRI}}$ & RE \\
    \midrule
    0.45 & 0.30 & 0.10 & 0.5 & 403 & 425 & 0.95 \\
     &  &  & 1 & 364 & 384 & 0.95 \\
     &  &  & 2 & 411 & 438 & 0.94 \\
     &  &  & 3 & 490 & 520 & 0.94 \\
    \addlinespace
    0.45 & 0.30 & 0.20 & 0.5 & 454 & 504 & 0.90 \\
     &  &  & 1 & 408 & 456 & 0.89 \\
     &  &  & 2 & 462 & 522 & 0.89 \\
     &  &  & 3 & 551 & 620 & 0.89 \\
    \addlinespace
    0.55 & 0.40 & 0.10 & 0.5 & 433 & 470 & 0.92 \\
     &  &  & 1 & 386 & 420 & 0.92 \\
     &  &  & 2 & 434 & 474 & 0.92 \\
     &  &  & 3 & 516 & 564 & 0.91 \\
    \addlinespace
    0.55 & 0.40 & 0.20 & 0.5 & 487 & 570 & 0.85 \\
     &  &  & 1 & 434 & 512 & 0.85 \\
     &  &  & 2 & 488 & 582 & 0.84 \\
     &  &  & 3 & 580 & 692 & 0.84 \\
    \bottomrule
  \end{tabular}
  \begin{tablenotes}[flushleft]\footnotesize
    \item $r = n_{1} / n_{0}$. RE $= N_{\mathrm{CC}} / N_{\mathrm{NRI}}$.
  \end{tablenotes}
  \end{threeparttable}
\end{table}
